% !TEX program = pdflatex
% !TEX encoding = UTF-8 Unicode
% Interim presentation, 25 Nov 2026 (DS500). Rules: docs/presentations.md.
% Every number comes from facts.tex next to this file.
\documentclass[aspectratio=169,10pt]{beamer}
\makeatletter\def\input@path{{../theme/}{./figures/}}\makeatother
\usetheme{reach}
% !TEX root = 2026-11-25-interim.tex
% Every number shown in the interim deck (25 Nov 2026) comes from here, with the run or document
% it was taken from. Update the whole block when the deck moves to a newer run; never type a
% number into the deck. Rule from docs/notes.md N13: always name profile and budget next to a share.

% ---- run used for results ----------------------------------------------------------------
% data/runs/20261110T223016Z/validate_results/report_results.md (status ok)
% This run has 7 categories (parks and playgrounds separate): it needs PR #57
% (feature/juju-categories, methodology v0.8). dev still has 6 ("Parks & playgrounds").
% Before presenting: #57 merged and a run from dev, or switch to a dev run with 6 categories.
\newcommand{\factRunId}{20261110T223016Z}
\newcommand{\factOsmDate}{8 Nov 2026}
\newcommand{\factCells}{1,699}
\newcommand{\factResidents}{90,355}
\newcommand{\factCategories}{7}
\newcommand{\factProfiles}{4}

% Share of residents per class, walking 3.39 km/h, 15 min (lack / insufficient / sufficient, %)
% Order: as in report_results.md. Used by the stacked bar chart of this deck.
\newcommand{\factLackPharmacyWalk}{43\,\%}   % 42.9 %

% ---- validation report (docs/validation_report.md, run 20261106T142630Z) -------------------
\newcommand{\factValRunId}{20261106T142630Z}
\newcommand{\factVOneOvercount}{18\,\%}      % V1: 585 polygon-only pairs = 18.3 % of matrix pairs, walking 3.39
\newcommand{\factVOneBarrier}{47\,\%}        % V1: polygon-only pairs across Neckar/rail/main road
\newcommand{\factVOneClass}{5\,\%}           % V1: cell x category classes that would differ
\newcommand{\factVTwoKita}{88\,\%}           % V2: childcare recall (87.5 %)
\newcommand{\factVTwoSchool}{95\,\%}         % V2: school recall
\newcommand{\factVThreeNear}{80--98\,\%}     % V3: picks within 500 m of a removed facility, mode \"created only\" (all deficits: 13--39 %)
\newcommand{\factVThreeRandom}{8--36\,\%}    % V3: same for random picks
\newcommand{\factVFourGap}{2.8\,\%}          % V4: max gap greedy vs. MILP (90 instances)
\newcommand{\factVFourMean}{0.17\,\%}        % V4: mean gap
\newcommand{\factVSixChange}{1.9\,\%}        % V6: residents x categories changing class, walking
\newcommand{\factVEightMoran}{0.72--0.88}    % V8: Moran's I on cells
\newcommand{\factVTenMedian}{0.5\,s}         % V10: median |delta| matrix vs. /route
\newcommand{\factVElevenChange}{3.3--3.6\,\%}% V11: residents x categories changing class

% ---- chart data (run \factRunId) -----------------------------------------------------------
% Share of residents per class [%], walking 3.39 km/h, 15 min; sorted by lack (ranking).
\pgfplotstableread{
category     lack  insufficient  sufficient
Pharmacy     42.9  0.0           57.1
Grocery      32.8  26.2          41.0
Parks        30.8  0.0           69.2
School       29.4  0.0           70.6
Doctors      17.5  23.1          59.4
Childcare    6.5   14.9          78.6
Playgrounds  4.2   0.0           95.8
}\factClassesWalk

% Share of residents in "lack" [%] per profile, 15 min (same order as above).
% Cycling 15 km/h: at most 2.6 % lack in every category (pharmacy).
\pgfplotstableread{
category     walkfive  walk  wheel
Pharmacy     25.7      42.9  45.3
Grocery      17.5      32.8  37.3
Parks        14.6      30.8  33.5
School       9.3       29.4  33.8
Doctors      8.3       17.5  22.1
Childcare    0.9       6.5   10.0
Playgrounds  0.8       4.2   7.1
}\factLackProfiles
\newcommand{\factLackCyclingMax}{3\,\%}  % 2.6 %, pharmacy

\input{isochrone}
\input{system}

\title[reach $\cdot$ Interim presentation]{reach: 15-minute accessibility for every resident of Tübingen}
\subtitle{Interim presentation $\cdot$ DS500 Data Science Project, Topic 9 ``The 15-Minute City''}
\author{Rico $\cdot$ Socke $\cdot$ Juju}
\institute{University of Tübingen, winter term 2026/27}
\date[25 Nov 2026]{25 November 2026}

\begin{document}

{\setbeamertemplate{headline}{}\setbeamertemplate{footline}{}
\begin{frame}
  \titlepage
\end{frame}}

% ============================================================================================
% Executive summary: the key message first (pyramid, workshop part 2). Speaker: TBD
\begin{frame}{reach shows who cannot reach everyday services within 15 minutes, and where a new
  facility would help most}
  \centering
  \begin{tikzpicture}[
      box/.style={draw=ink,line width=0.6pt,text width=40mm,align=flush left,inner sep=2.5mm,
        font=\footnotesize,minimum height=27mm,anchor=north},
      head/.style={font=\footnotesize\bfseries,text=accent}]
    \node[fill=ink,text=white,text width=132mm,align=center,inner sep=2.5mm,font=\small\bfseries]
      (key) at (0,0)
      {Accessibility in Tübingen is uneven: it depends on where you live and how you travel};
    \node[box] (a) at (-46mm,-11mm) {{\usebeamercolor[fg]{alerted text}\bfseries Measure}\\[1mm]
      Every populated 100\,m cell, \factCategories{} categories, \factProfiles{} ways of
      travelling, along the street network};
    \node[box] (b) at (0,-11mm) {{\usebeamercolor[fg]{alerted text}\bfseries Find}\\[1mm]
      \factLackPharmacyWalk{} of residents cannot walk to a pharmacy within 15 minutes
      (3.39\,km/h)};
    \node[box] (c) at (46mm,-11mm) {{\usebeamercolor[fg]{alerted text}\bfseries Recommend}\\[1mm]
      Zones for one new facility per category, ranked by unmet demand removed; within
      \factVFourGap{} of the optimum};
    \foreach \n in {a,b,c} \draw[ink,line width=0.6pt] (key.south) -- (\n.north);
    \node[font=\footnotesize,text=muted,anchor=north] at (0,-41mm)
      {Automated from open data with one command; 9 of 11 validation experiments done};
  \end{tikzpicture}
  \source{run \factRunId{} (OSM \factOsmDate); validation report, run \factValRunId.}
\end{frame}

% ============================================================================================
\section{Problem}
% Speaker: TBD

\begin{frame}{The 15-minute city asks whether residents reach daily needs without a car; we
  make that question measurable and actionable for Tübingen}
  \begin{columns}[T,onlytextwidth]
    \column{0.47\textwidth}
      \textbf{The idea}\par\smallskip
      \begin{itemize}
        \item Everyday services within a short walk or ride from home
              (Moreno et al., 2021)
        \item Mouratidis (2026): count \emph{disaggregated} facilities per category and classify
              access as lack, insufficient or sufficient
        \item Applied once, largely by hand, to Copenhagen
      \end{itemize}
    \column{0.47\textwidth}
      \textbf{What we add}\par\smallskip
      \begin{itemize}
        \item \textbf{Automation:} open data in, validated results out, with one command
        \item \textbf{Tübingen} with German open data (OSM, NVBW GTFS, Zensus 2022)
        \item \textbf{Inclusive profiles:} slow walking, wheelchair, cycling
        \item \textbf{A recommender:} where would one new facility help most?
      \end{itemize}
  \end{columns}
  \source{Moreno et al. (2021); Mouratidis (2026), full references in the backup.}
\end{frame}

\begin{frame}{An isochrone shows every place you can reach from one point within a fixed time}
  \begin{columns}[T,onlytextwidth]
    \column{0.36\textwidth}
      \begin{block}{Isochrone}
        \centering $I(o,t) = \{\,x \;:\; \tau_m(o,x) \le t\,\}$\par\smallskip
        \raggedright\footnotesize $o$: origin, $t$: time budget, $\tau_m$: travel time
        along the street network for profile $m$
      \end{block}
      \smallskip
      \begin{itemize}\small
        \item Our budget: $t = 15$ min
        \item At 3.39\,km/h: $\approx 850$\,m \emph{along streets}
        \item Follows the network, stops at barriers like the Neckar
      \end{itemize}
    \column{0.62\textwidth}
      \centering\isochronedefinition\par
      {\tiny\color{muted}Schematic. Filled squares: reachable within 15 min; contours at 5, 10
      and 15 min.}
  \end{columns}
  \source{methodology.md §1 (profiles); isochrones are computed by Valhalla.}
\end{frame}

% ============================================================================================
\section{Method}
% Speaker: TBD

\begin{frame}{We measure all \factCells{} populated 100\,m cells (\factResidents{} residents)
  for seven categories and four ways of travelling}
  \begin{columns}[T,onlytextwidth]
    \column{0.31\textwidth}
      \textbf{Who: Zensus 2022}\par\smallskip
      \footnotesize Residents per 100\,m grid cell; the cell centroid is the origin\par
      \bigskip\normalsize
      \textbf{What: OpenStreetMap}\par\smallskip
      \footnotesize Street network and facilities: grocery, pharmacy, doctors, childcare,
      school, parks, playgrounds\par
      \bigskip\normalsize
      \textbf{Transit: NVBW GTFS}\par\smallskip
      \footnotesize Weekday departures at stops within a 10-minute walk
    \column{0.63\textwidth}
      \centering
      \footnotesize
      \begin{tabular}{@{}lrr@{}}
        \toprule
        Profile & Speed & Reach in 15 min \\
        \midrule
        Walking (inclusive default) & 3.39 km/h & 850 m \\
        Walking (comparison) & 5.0 km/h & 1,250 m \\
        Wheelchair & 3.39 km/h & 850 m \\
        Cycling & 15 km/h & 3,750 m \\
        \bottomrule
      \end{tabular}\par\medskip
      \raggedright
      Reach = straight-line limit; the network time decides. Slopes slow walkers and
      wheelchairs down (Valhalla's slope model).
  \end{columns}
  \source{run \factRunId{} (OSM \factOsmDate); methodology.md §1--§3.}
\end{frame}

\begin{frame}{We count facilities along the network, because isochrone polygons would add
  \factVOneOvercount{} unreachable cell--facility pairs}
  \begin{columns}[T,onlytextwidth]
    \column{0.57\textwidth}
      \centering\isochronetrap
    \column{0.40\textwidth}
      \textbf{Matrix, not polygon:} for every cell, one Valhalla time matrix to all candidate
      facilities; a facility counts if its travel time $\le 15$ min.\par\bigskip
      \textbf{Test V1} (150 cells, walking 3.39\,km/h):
      \begin{itemize}
        \item Polygons add \factVOneOvercount{} unreachable pairs and miss almost none
        \item \factVOneBarrier{} of these lie across the Neckar, rail lines or main roads
        \item \factVOneClass{} of classes would change
      \end{itemize}
  \end{columns}
  \source{validation report V1, run \factValRunId; architecture.md §6.3 (decision D2).}
\end{frame}

\begin{frame}{Each cell gets a class per category, and a deficit that weighs missing access by
  residents}
  \begin{columns}[T,onlytextwidth]
    \column{0.48\textwidth}
      $A_{i,c}$ = facilities of category $c$ reachable from cell $i$ in 15 min,
      $k_c$ = facilities needed (grocery, doctors, childcare: 2; others: 1)\par\bigskip
      \begin{tabular}{@{}ll@{}}
        \classchip{lack}{Lack} & $A_{i,c} = 0$ \\[2mm]
        \classchip{insufficient}{Insufficient} & $0 < A_{i,c} < k_c$ \\[2mm]
        \classchip[ink]{sufficient}{Sufficient} & $A_{i,c} \ge k_c$ \\
      \end{tabular}
    \column{0.48\textwidth}
      \textbf{Deficit} (unmet demand, residents $P_i$)
      \[ U_{i,c} = P_i \cdot \max\!\Bigl(0,\; 1 - \frac{A_{i,c}}{k_c}\Bigr) \]
      \footnotesize Example: 120 residents, one supermarket ($k = 2$) $\Rightarrow U = 60$.\par
      \bigskip\normalsize
      We store only raw counts and times; classes and deficits are derived when read, so a new
      $k_c$ needs no re-routing.
  \end{columns}
  \source{methodology.md §2, §4.}
\end{frame}

% ============================================================================================
\section{Results}
% Speaker: TBD

\begin{frame}{More than four in ten residents cannot walk to a pharmacy within 15 minutes}
  \centering
  \begin{tikzpicture}
    \begin{axis}[
        reach ranking, reach box legend, xbar stacked, bar width=4.2mm,
        yticklabels from table={\factClassesWalk}{category}, xmin=0, xmax=100,
        xlabel={Share of residents [\%], walking 3.39\,km/h, 15 min},
      ]
      \addplot[fill=lack,draw=none,
        nodes near coords={\pgfmathfloattofixed{\pgfplotspointmeta}\ifdim\pgfmathresult pt<5pt\else
          \pgfmathprintnumber[fixed,precision=0]{\pgfplotspointmeta}\fi},
        nodes near coords style={font=\scriptsize\bfseries,text=white,anchor=east,
          xshift=-0.3mm},
        point meta=explicit]
        table[x=lack,y expr=\coordindex,meta=lack] {\factClassesWalk};
      \addplot[fill=insufficient,draw=none]
        table[x=insufficient,y expr=\coordindex] {\factClassesWalk};
      \addplot[fill=sufficient,draw=none]
        table[x=sufficient,y expr=\coordindex] {\factClassesWalk};
      \legend{Lack, Insufficient, Sufficient}
    \end{axis}
  \end{tikzpicture}
  \source{run \factRunId{} (OSM \factOsmDate), \factResidents{} residents in \factCells{} cells.
    Insufficient exists only where $k_c = 2$.}
\end{frame}

\begin{frame}{Faster walkers reach far more; pharmacy and grocery stay the weakest on foot and
  by wheelchair}
  \centering
  \begin{tikzpicture}
    \begin{axis}[
        reach ranking, reach legend,
        yticklabels from table={\factLackProfiles}{category}, xmin=0, xmax=50,
        xmajorgrids, grid style={rule,line width=0.3pt},
        xlabel={Share of residents in lack [\%], 15 min},
      ]
      \addplot[draw=none,mark=none,forget plot,
        error bars/.cd, x dir=plus, x explicit, error mark=none,
        error bar style={line width=2.2pt,rule}]
        table[x=walkfive,y expr=\coordindex,
          x error expr=\thisrow{wheel}-\thisrow{walkfive}] {\factLackProfiles};
      \addplot[only marks,mark=*,mark size=2.6pt,mark options={fill=white,draw=ink}]
        table[x=walkfive,y expr=\coordindex] {\factLackProfiles};
      \addplot[only marks,mark=*,mark size=2.6pt,mark options={fill=ink,draw=ink}]
        table[x=walk,y expr=\coordindex] {\factLackProfiles};
      \addplot[only marks,mark=diamond*,mark size=3.4pt,mark options={fill=accent,draw=accent}]
        table[x=wheel,y expr=\coordindex] {\factLackProfiles};
      \legend{Walking 5.0\,km/h, Walking 3.39\,km/h, Wheelchair 3.39\,km/h}
    \end{axis}
  \end{tikzpicture}
  \source{run \factRunId{} (OSM \factOsmDate). Cycling 15\,km/h: at most \factLackCyclingMax{}
    lack in every category. District rankings stay similar under sensitivity changes (V5).}
\end{frame}

\begin{frame}{A greedy recommender picks the zone that removes the most unmet demand, close to
  the exact optimum}
  \begin{columns}[T,onlytextwidth]
    \column{0.58\textwidth}
      \textbf{Gain} of a new facility of category $c$ at zone $x$:
      \[ \Delta W(x,c) = \sum_{i \in C(x)} P_i \left[\min\left(1,\frac{A_{i,c}+1}{k_c}\right) -
         \min\left(1,\frac{A_{i,c}}{k_c}\right)\right] \]\vspace{-2mm}
      \begin{enumerate}\small
        \item Candidates: 100\,m zones in the built-up area
        \item Pick the largest $\Delta W$, update the counts, repeat
        \item Submodular objective: greedy reaches $\ge 1 - 1/e$ of the optimum
      \end{enumerate}
    \column{0.36\textwidth}
      \keyfigure{$\le$\,\factVFourGap}{below the exact MILP optimum\\(90 instances, mean
        \factVFourMean)}\par\bigskip
      \keyfigure{\factVThreeNear}{of picks near a removed facility when only the removed
        deficit counts (back-test; random: \factVThreeRandom)}\par\bigskip
      {\scriptsize\color{muted} Accessibility only: zoning, parcels and capacity are out of
      scope; results name zones, not sites.}
  \end{columns}
  \source{methodology.md §5; validation report V3, V4, run \factValRunId.}
\end{frame}

% ============================================================================================
\section{System}
% Speaker: TBD

\begin{frame}{One command turns open data into validated, published results}
  \vspace{4mm}
  \centering\pipelinefigure\par
  \vspace{6mm}
  \begin{columns}[T,onlytextwidth]
    \column{0.31\textwidth}\footnotesize\textbf{Reproducible:} every run keeps its inputs,
      SHA-256 hashes and a manifest per stage
    \column{0.31\textwidth}\footnotesize\textbf{One routing engine:} Valhalla for counting,
      isochrones and live queries in the app
    \column{0.31\textwidth}\footnotesize\textbf{Fast:} the full count for Tübingen takes
      about 3.5 min on a laptop
  \end{columns}
  \source{architecture.md §6--§7; \texttt{make refresh}.}
\end{frame}

\begin{frame}{Eleven experiments test the method, and the main design choices hold}
  \vspace{4mm}
  \begin{columns}[T,onlytextwidth]
    \column{0.24\textwidth}\keyfigure{\factVOneOvercount}{V1: over-count avoided by matrix
      counting}
    \column{0.24\textwidth}\keyfigure{\factVTwoSchool}{V2: OSM recall for schools
      (childcare \factVTwoKita)}
    \column{0.24\textwidth}\keyfigure{\factVSixChange}{V6: class changes with building-based
      origins instead of centroids}
    \column{0.24\textwidth}\keyfigure{\factVTenMedian}{V10: median difference matrix vs.\
      single routes}
  \end{columns}
  \vspace{8mm}
  \begin{itemize}
    \item Deficits are spatially clustered (V8, Moran's I \factVEightMoran{} on cells)
    \item Slopes change \factVElevenChange{} of classes, mostly on the hillsides (V11)
    \item Still open: plausibility with local knowledge (V7), reproducibility by a teammate (V9)
  \end{itemize}
  \source{validation report V1--V11, run \factValRunId.}
\end{frame}

\begin{frame}{Live demo: anyone can click a place and see what is missing within 15 minutes}
  \begin{columns}[T,onlytextwidth]
    \column{0.55\textwidth}
      \begin{enumerate}
        \item Class map: pharmacy, walking 3.39\,km/h
        \item Pin drop in an underserved area: what is reachable, what is missing, how long
        \item Recommended zone: which residents would gain access
      \end{enumerate}
    \column{0.40\textwidth}
      \begin{block}{If the demo fails}
        \footnotesize We show screenshots of the same three steps instead.
      \end{block}
  \end{columns}
\end{frame}

% ============================================================================================
\section{Next steps}
% Speaker: TBD

\begin{frame}{By the submission on 22 January we close the open validations and host the app}
  \begin{columns}[T,onlytextwidth]
    \column{0.48\textwidth}
      \textbf{Validation}
      \begin{itemize}
        \item V7: local-knowledge check of the worst areas
        \item V9: a teammate reruns everything from the README
        \item Map the missing Kitas and schools in OSM (V2)
      \end{itemize}
    \column{0.48\textwidth}
      \textbf{Product}
      \begin{itemize}
        \item Hosted app and read-only database access for supervisors
        \item Weekly automatic refresh on the server
        \item Calibrate the transit supply classes to Tübingen
      \end{itemize}
  \end{columns}
  \vspace{6mm}
  {\footnotesize\color{muted}Known limit: the wheelchair profile avoids steps but not steep
  ramps, so wheelchair access on hillsides is optimistic.}
  \source{docs/notes.md N11, N13; validation report.}
\end{frame}

{\setbeamertemplate{headline}{}
\begin{frame}
  \vfill
  {\Large\bfseries Thank you. Questions?\par}
  \vskip4mm
  {\color{accent}\rule{25mm}{2pt}}\par\vskip4mm
  Accessibility in Tübingen is uneven: it depends on where you live and how you travel.
  reach measures it for every resident and shows where a new facility helps most.\par
  \vfill
\end{frame}}

% ============================================================================================
\backupslides
\section*{Backup}

\begin{frame}{Backup: references and data sources}
  \footnotesize
  \begin{itemize}
    \item Moreno, C., Allam, Z., Chabaud, D., Gall, C., \& Pratlong, F. (2021). Introducing the
          ``15-Minute City'': Sustainability, resilience and place identity in future
          post-pandemic cities. \emph{Smart Cities}, 4(1), 93--111.
    \item Mouratidis, K. (2026). \emph{Travel Behaviour and Society}, 42, 101151.
  \end{itemize}
  \medskip
  \begin{tabular}{@{}lll@{}}
    \toprule
    Source & Used for & Licence \\
    \midrule
    OpenStreetMap (Geofabrik extract) & Streets, facilities & ODbL 1.0 \\
    NVBW GTFS (naldo) & Stops, departures & dl-de/by-2-0, ODbL \\
    Destatis Zensus 2022, 100\,m grid & Residents per cell & dl-de/by-2-0 \\
    Valhalla (open source) & Routing, matrix, isochrones & MIT \\
    \bottomrule
  \end{tabular}
  \source{README.md (data sources); Valhalla: licence file of the project repository.}
\end{frame}

\begin{frame}{Backup: validation experiments V1--V11}
  \scriptsize
  \begin{tabular}{@{}llp{95mm}@{}}
    \toprule
    \# & Experiment & Key result \\
    \midrule
    V1 & Matrix vs.\ isochrone polygons & Polygons add 18--20\,\% false pairs, 47\,\% across barriers \\
    V2 & OSM completeness & Childcare recall 88\,\%, precision 98\,\%; schools 95\,\% / 95\,\% \\
    V3 & Recommender back-test & 80--98\,\% of picks within 500\,m of a removed facility (only the removed deficit counts) \\
    V4 & Greedy vs.\ exact MILP & At most 2.8\,\% below the optimum (mean 0.17\,\%) \\
    V5 & Sensitivity ($k$, speed, budget) & Shares move strongly; district rankings stable
      (median Spearman 0.78--0.95) \\
    V6 & Centroid vs.\ building origin & 1.9\,\% (walking) / 2.8\,\% (wheelchair) change class \\
    V7 & Local knowledge & Not yet run \\
    V8 & Spatial statistics & Moran's I 0.72--0.88 on cells; none between districts \\
    V9 & Reproducibility & Not yet run \\
    V10 & Matrix vs.\ single routes & Median $|\Delta|$ 0.5\,s; 0 of 900 false negatives \\
    V11 & With vs.\ without slopes & 3.3--3.6\,\% change class, mostly on hillsides \\
    \bottomrule
  \end{tabular}
  \source{docs/validation\_report.md, run \factValRunId.}
\end{frame}

\begin{frame}{Backup: architecture at runtime}
  \vspace{10mm}
  \centering\architecturefigure\par
  \vspace{10mm}
  \footnotesize Docker Compose on one small server; Caddy serves the web app over HTTPS. The API
  reads only the \texttt{current} views of the last validated run.
  \source{architecture.md §3.}
\end{frame}

\end{document}
